\section{Limitations}\label{sec:limits}
% src: S-PAPER-1, S-PAPER-2, S-PAPER-3, S-PAPER-6, S-PAPER-10
Some markets have no usable winner label, and we drop them.
In our full run this is \NNoOutcome{} of \NMarketsInScope{} markets (\PctNoOutcome).
Our rule is wider than the dataset paper's.
It also drops markets whose outcome-1 label never traded.
The paper says only a small subset lacks a winner.

Markets keep trading after \texttt{close\_at}.
We avoid this by pricing one day before the scheduled close.
But the choice rests on a one-month pilot.

Sports is the largest topic in scope and in the kept markets.
It is \SportsShareInScope{} of the \NMarketsInScope{} markets in scope and \SportsShareKept{} of the kept markets.
Its survival is \CatSportsSurvival{} against \PctKept{} overall.
So Sports is slightly under-represented against its scope.
Sports loses markets mainly to no price (\CatSportsDropNoPriceN) or a stale price (\CatSportsDropStalePriceN).
Both are measured \DaysShort{} day before the close.
It loses almost none to the Up/Down rule (\CatSportsDropUpdownN).
The dataset paper notes that Sports markets fall disproportionately into Neg-Risk or Up/Down \citep{qin2026polymarketv1}.
Neg-Risk markets are outside the layer we use.

Our selection is narrow.
Only \NMarketsKept{} of \NMarketsInScope{} markets pass the filters.
\NNoPrice{} markets have no price \DaysShort{} day before the close (\PctNoPrice).
The minimum-trade filter is a choice, and we show the result without it.
\Cref{sec:app-m4} compares kept and dropped markets (M4).
Markets on one event move together.
The binomial and chi-square tests assume independent markets, so the p-values are optimistic.
The z intervals use the same assumption, so they may be too narrow.
The role tests compare groups that share addresses and trades.
No event key exists in the data, so we cannot size this dependence exactly.
Two proxy keys address it in part (\Cref{sec:app-m3}).
The proxies are not true event keys.
With this many rows, a tiny gap can give a tiny p-value.
So we read the effect sizes first.
The intervals do not escape the independence assumption.

The k-means result is exploratory.
The clusters describe the sample.
A wallet is an address, and one owner may control many addresses.
Volume concentration is not the same as the power to set prices.

We changed one pre-registered rule after seeing a full Q1 result.
This is the tie rule in \Cref{sec:methods-prereg}.
We report it as an amendment.
The amendment records that the verdict did not change.

\Cref{tab:checks} sums up the robustness checks R1 and R2.
